% ==============================================================================
% SECCIÓN 4: ORGANIZACIÓN DEL EQUIPO (ROLES Y NORMAS)
% ==============================================================================
\section{Organización del Equipo de Trabajo (Estructura, Normas)}

\subsection{Asignación Nominal de Roles Iniciales}
Se seguirá una división de roles entre los 6 integrantes del equipo, rotando aquellos roles menos importantes en cada iteración para repartir dicha carga. La asignación inicial de roles del equipo es:
\begin{itemize}[leftmargin=*]
  \item \textbf{Pablo Hernández Ibáñez:} Moderador (Rotativo)
  \item \textbf{Yeray Rodríguez Navas:} Gestor de Calidad (\textbf{Fijo})
  \item \textbf{Pablo de la Torre Roldán:} Catalogador (Rotativo)
  \item \textbf{José Rodríguez Fernández:} Presentador (Rotativo)
  \item \textbf{Julián Carrión Tovar:} Gestor de Usabilidad y Accesibilidad (\textbf{Fijo})
  \item \textbf{Miguel Ángel Luque:} Coordinador (Rotativo)
\end{itemize}

\subsection{Responsabilidades de Cada Rol}
Cada rol tendrá las siguientes responsabilidades:
\begin{itemize}[leftmargin=*]
  \item \textbf{Coordinador:} Organiza el trabajo del equipo, distribuye tareas, controla los plazos y realiza el seguimiento general del proyecto. Se encarga de avisar sobre las próximas entregas y de planificar el tiempo de trabajo correspondiente a cada iteración. Actúa como portavoz de gestión ante el profesorado y supervisa el equilibrio de la dedicación horaria.
  \item \textbf{Catalogador:} Responsable de recopilar, analizar y clasificar la información generada por el grupo en las distintas tareas asociadas a las prácticas. Se encarga de revisar y consolidar la información añadida a la documentación en cada iteración y, cuando sea posible, de revisar la documentación de la iteración anterior en coordinación con el Gestor de Calidad. Asimismo, custodia y organiza la documentación en el repositorio institucional (\texttt{docs/}) y coordina la publicación de actas de reunión.
  \item \textbf{Moderador:} Responsable de plantear y moderar los debates, así como de seleccionar las ideas y decisiones grupales. En caso de indecisión, tiene la última palabra tras consultar con el Gestor de Calidad sobre la opción más adecuada. Asimismo, se encarga de gestionar y documentar los cambios en los documentos de especificación.
  \item \textbf{Presentador:} Se encarga de realizar todas las comunicaciones con el cliente, así como de elaborar presentaciones que muestren el avance y los resultados del proyecto. Prepara los resúmenes que se comunicarán al cliente o profesor y decide, junto con los roles pertinentes, qué contenidos y preguntas se expondrán en cada iteración.
  \item \textbf{Gestor de la Calidad (Rol Fijo):} Responsable de asegurar la calidad de los productos generados y del proceso de desarrollo. Vela por que la documentación y las técnicas aplicadas cumplan con los estándares definidos:
    \begin{itemize}
      \item Asegurarse de que se usan herramientas estandarizadas para especificación y diseño (diagramas UML en Visual Paradigm y plantillas de ingeniería de software).
      \item Asegurarse de que se organiza bien el trabajo: celebración regular de reuniones, toma de actas formales, reparto equitativo de tareas y uso de recursos hardware y software correctos.
      \item Revisar o asegurarse de que se revisen y prueben todos los entregables (verificando correctitud y completitud) y de que se finalizan a tiempo o se planifica adecuadamente cualquier cambio de fecha y sus consecuencias.
      \item Supervisar los criterios de aceptación, revisiones de código, análisis estático continuo y auditorías periódicas de calidad del proceso y producto.
    \end{itemize}
  \item \textbf{Gestor de Usabilidad y Accesibilidad (Rol Fijo):} Responsable de planificar, coordinar y velar por la usabilidad y la accesibilidad del proyecto. Se encarga de investigar e informar al equipo sobre las buenas prácticas, técnicas y herramientas a utilizar en estas materias. Asimismo, define los requisitos de accesibilidad, formula las preguntas necesarias de las consultas que se harán y verifica su correcto cumplimiento mediante pruebas de usuario, testing o entrevistas. Para el desarrollo e implementación de estas actividades contará con el apoyo de los roles correspondientes:
    \begin{itemize}
      \item Formarse y formar a sus compañeros en guías de usabilidad y accesibilidad web y móvil (WCAG~2.1~AA, heurísticas de accesibilidad cognitiva y diseño accesible).
      \item Planificar qué guías considerar en función de los usuarios finales del proyecto (alumnado PTVAL con NEAE, docentes y familias).
      \item Formarse en heurísticas de accesibilidad y emplearlas rigurosamente para evaluar y probar la aplicación.
      \item Aplicar lectores de pantalla (\textbf{TalkBack} en Android y \textbf{VoiceOver} en iOS) para validar la accesibilidad móvil.
      \item Aplicar validadores de accesibilidad (\textbf{Lighthouse}, \textbf{WAVE}, Google Accessibility Scanner) para evaluar contrastes cromáticos, etiquetas y comprensión.
      \item Elaborar informes periódicos del progreso del trabajo en cuanto a usabilidad y accesibilidad.
    \end{itemize}
\end{itemize}

\subsection{Fijeza, Esquema de Rotación y Protocolo de Ausencias}
Los roles de \textbf{Gestor de Calidad} y \textbf{Gestor de Usabilidad y Accesibilidad} serán \textbf{fijos} a los responsables correspondientes (Yeray Rodríguez Navas y Julián Carrión Tovar), mientras que el resto rotará en cada iteración de trabajo.

\begin{infobox}[title=Regla de Rotación Cíclica]
La secuencia obligatoria de rotación en cada cambio de iteración es:\\
\centerline{\large $\dots \longrightarrow \textbf{Coordinador} \longrightarrow \textbf{Catalogador} \longrightarrow \textbf{Moderador} \longrightarrow \textbf{Presentador} \longrightarrow \dots$}
\end{infobox}

\noindent\textbf{Protocolo de Suplencias y Ausencias:}
\begin{enumerate}[leftmargin=*]
  \item En caso de que algún integrante con cualquiera de los roles de Gestor (Calidad o Accesibilidad) no pueda acudir a una iteración de trabajo, dicho rol será cedido temporalmente al que en dicha reiteración tenga el rol de \textbf{Coordinador} (y al \textbf{Presentador} en caso de inasistencia de ambos).
  \item Si algún rol rotativo no está presente en una iteración, se le cederá durante esa iteración al responsable de ese rol en la siguiente iteración.
\end{enumerate}

\subsection{Matriz de Rotación de Roles por Iteración}

\vspace{0.2cm}
\noindent
{\small
\begin{tabularx}{\textwidth}{|>{\raggedright\arraybackslash\bfseries}p{4.2cm}|>{\centering\arraybackslash}p{2.3cm}|>{\centering\arraybackslash}X|>{\centering\arraybackslash}X|>{\centering\arraybackslash}X|}
\hline
\rowcolor{tableHeader}
Integrante & Fase Inicial & Iteración 1 & Iteración 2 & Iteración 3 \\
\hline
Miguel Ángel Luque & \textbf{Coordinador} & Catalogador & Moderador & Presentador \\
\hline
Pablo de la Torre Roldán & \textbf{Catalogador} & Moderador & Presentador & Coordinador \\
\hline
Pablo Hernández Ibáñez & \textbf{Moderador} & Presentador & Coordinador & Catalogador \\
\hline
José Rodríguez Fernández & \textbf{Presentador} & Coordinador & Catalogador & Moderador \\
\hline
Yeray Rodríguez Navas & \textbf{Gestor Calidad} & \textbf{Gestor Calidad} & \textbf{Gestor Calidad} & \textbf{Gestor Calidad} \\
\hline
Julián Carrión Tovar & \textbf{Gestor Accesib.} & \textbf{Gestor Accesib.} & \textbf{Gestor Accesib.} & \textbf{Gestor Accesib.} \\
\hline
\end{tabularx}
}
\vspace{0.3cm}

\subsection{Normas de Convivencia y Compromiso Operativo}
\begin{enumerate}[leftmargin=*]
  \item \textbf{Puntualidad:} Margen máximo de cortesía de 5 minutos en reuniones presenciales y virtuales.
  \item \textbf{Notificación Previa de Bloqueos:} Notificar obligatoriamente cualquier impedimento o retraso con al menos \textbf{48 horas de antelación}.
  \item \textbf{Registro Honesto de Horas:} Cumplimentación semanal estricta de la ficha de dedicación horaria individual.
\end{enumerate}
